\chapter{Literature Review}
\label{ch:literature}

This chapter situates the thesis within three intersecting literatures: (i) the empirical evaluation of clean-energy policy, with a focus on hydrogen; (ii) the real-options theory of irreversible investment; and (iii) recent advances in causal inference, particularly modern difference-in-differences and honest sensitivity analysis.

\section{Empirical evaluation of hydrogen and clean-energy policy}
\label{sec:lit_hydrogen}

\subsection{The hydrogen implementation gap}

The empirical literature on hydrogen-project realisation is dominated by descriptive tracking of capacity announcements versus delivered FIDs. \cite{odenweller2025green} construct what they term the \emph{ambition gap} (the difference between announced 2030 capacity and 1.5\textdegree C-consistent demand) and the \emph{implementation gap} (the difference between announced and realised capacity at scheduled launch dates). Tracking 190 projects over three years, they document a 2023 implementation gap of approximately 93\%: only 7\% of capacity announcements were on schedule. Their projection requires global subsidies of US\$1.3 trillion (range: \$0.8--2.6 trillion) to realise the announced project pipeline, far exceeding announced subsidies.

\cite{ueckerdt2024cost} provide a complementary cost-competitiveness analysis. Under current natural-gas prices, green hydrogen is initially more than eight times as expensive as natural gas; even under an ambitious carbon-price pathway, it remains 3.5 times more expensive in 2024, with cost parity not reached before 2045. These cost gaps imply that voluntary uptake without policy support is unlikely.

\cite{hydrogen2024insights} and the IEA \cite{iea2024hydrogen} provide industry-side tracking with broader capacity coverage but less methodological rigour. The S\&P Global Commodity Insights team maintains the dataset used in this thesis, including project-level status updates that have not been previously exploited for causal evaluation.

\subsection{Policy-evaluation studies}

Direct empirical evaluation of specific hydrogen-policy mechanisms remains limited. \cite{neuwirth2023industry} simulate the EU Innovation Fund's projected impact on industrial hydrogen adoption but rely on bottom-up cost modelling rather than ex-post causal estimates. \cite{ricks2023minimizing} analyse the additionality and temporal-matching requirements of the US 45V production credit but stop short of project-survival outcomes. The closest predecessors to the present analysis are \cite{hosseini2022real} and \cite{lloyd2022maritime}, who apply real-options frameworks to specific electrolysis and maritime-shipping cases but do not exploit cross-jurisdictional variation.

To the author's knowledge, no prior study comparatively evaluates the four carrot-policy archetypes considered here within a unified causal framework.

\subsection{Recent developments (2024--2026): the cancellation wave}

The empirical literature on hydrogen-project realisation has been overtaken by industry events since 2024. The \cite{iea2025hydrogen} \emph{Global Hydrogen Review 2025} reports that the announced 2030 production capacity declined for the first time since the database's inception: from 49 Mtpa in the 2024 review to 37 Mtpa in 2025, a decline of nearly one-quarter in a single calendar year. Electrolysis projects accounted for more than 80\% of the total drop. The IEA attributes the recalibration to high production costs (renewable hydrogen one-and-a-half to six times more costly than unabated fossil-based production), uncertain demand signals, financing hurdles, delays to incentive disbursement (notably under US 45V), regulatory uncertainty, and operational complications. 

\cite{odenweller2025green} document the same pattern from a different methodological angle, finding that the historical realisation rate of announced low-emissions hydrogen projects fell from approximately 7\% historically to 4\% by 2024. The implication is that the industry has entered a regime change, and contemporaneous empirical work on hydrogen policy must accommodate this structural break rather than treating the 2020--2023 announcement era as representative of long-run dynamics.

The present thesis is positioned to contribute to this evolving literature by providing the first systematic econometric quantification of which policy and contracting mechanisms differentially affect project survival, conditional on this regime change. The S\&P Global database snapshot of March 2026 used here captures the cancellation wave's complete cross-section; the carrot-policy and offtake-effect estimates of Chapters \ref{ch:results_did} and \ref{ch:results_offtake} are therefore identified on the most policy-relevant variation period currently observable.

\subsection{Adjacent literature: solar and wind policy evaluation}

The broader clean-energy policy-evaluation literature provides useful methodological precedent. \cite{bento2018solar} use a panel difference-in-differences design to identify solar-subsidy effects on installed capacity, finding heterogeneous effects across US states. \cite{aldy2016environment} review feed-in-tariff versus renewable-portfolio-standard designs, concluding that mechanism design matters as much as subsidy intensity --- a finding that motivates the present focus on \emph{mechanism-type} rather than \emph{subsidy-magnitude} comparison. \cite{popp2010innovation} examine the patenting response to clean-energy R\&D subsidies, finding pronounced country-specific heterogeneity that anticipates our cross-jurisdictional approach.

\section{Real options and irreversible investment}
\label{sec:lit_realoptions}

\subsection{The Dixit-Pindyck framework}

The canonical framework for analysing irreversible investment under uncertainty is \cite{dixit1994investment}. A project of present value $V$ (stochastic) and sunk cost $I$ (fixed) has positive option value to delay investment whenever the underlying $V_t$ follows geometric Brownian motion with non-zero volatility $\sigma$. The optimal exercise (FID) rule is to invest when $V \geq V^*$, where

\begin{equation}
\frac{V^*}{I} = \frac{\beta_1}{\beta_1 - 1}, \quad
\beta_1 = \frac{1}{2} - \frac{r-\delta}{\sigma^2}
       + \sqrt{\left(\frac{r-\delta}{\sigma^2} - \frac{1}{2}\right)^2 + \frac{2r}{\sigma^2}}
\label{eq:lit_dixit_pindyck}
\end{equation}

with $r$ the risk-free rate and $\delta$ the convenience-yield-like dividend. Higher $\sigma$ raises $V^*/I$ monotonically: greater uncertainty makes waiting more valuable, suppressing investment.

\subsection{Applications to energy investment}

\cite{pindyck1991irreversibility} provides the foundational application of real-options theory to energy. \cite{fuss2012impact} apply the framework to carbon-pricing uncertainty in power-sector investment, showing that policy-volatility itself can suppress investment even when expected returns are positive. \cite{barradale2014investing} document that wind-investment cycles are driven primarily by policy-uncertainty timing rather than by underlying renewable economics.

For hydrogen specifically, \cite{hosseini2022real} apply a binomial-tree real-options model to electrolyser projects, finding that early subsidies can shift optimal-exercise timing forward by 2--4 years. \cite{lloyd2022maritime} extend the framework to maritime ammonia, demonstrating that demand-aggregation can functionally reduce $\sigma$.

Our theoretical contribution (Chapter \ref{ch:theory}) extends this literature by formalising the mechanism-design implication: different policy types attack different parameters of the Dixit-Pindyck framework, and their relative effectiveness depends on the sector-specific $(V/I, \sigma)$ distribution.

\section{Causal inference: difference-in-differences and beyond}
\label{sec:lit_did}

\subsection{The two-way fixed effects (TWFE) problem}

The classical difference-in-differences (DiD) estimator with two-way fixed effects \cite{ashenfelter1985using} has been the workhorse of policy-evaluation research for four decades. However, recent literature has documented serious identification problems when treatment timing is staggered and effects are heterogeneous across cohorts \cite{goodman2021difference,dechaisemartin2020two,callaway2021difference}. Under heterogeneous effects, TWFE can produce coefficients with the wrong sign relative to any individual treatment effect.

\subsection{Modern DiD estimators}

\cite{sun2021estimating} develop an interaction-weighted estimator that uses only never-treated or last-treated units as controls, eliminating the "forbidden comparison" problem of staggered TWFE. \cite{callaway2021difference} (CS) propose group-time average treatment effects $ATT(g,t)$ that can be aggregated under various weighting schemes. \cite{borusyak2024revisiting} (BJS) develop an imputation estimator that, under a parallel-trends assumption on untreated potential outcomes, achieves efficiency bounds and is robust to heterogeneous effects. \cite{dechaisemartin2024intertemporal} characterise intertemporal treatment effects in heterogeneous-treatment settings, providing identification guarantees under weaker conditions than CS or BJS.

\cite{roth2023trending} synthesise these developments in the most cited recent review of the DiD literature, classifying estimators by their relaxations of canonical assumptions (multiple periods, heterogeneous timing, parallel-trends violations, alternative inference). The recommendations from this synthesis directly motivate the present thesis's triangulation strategy: rather than committing to a single estimator, robust empirical claims should rest on agreement across multiple specifications that relax different canonical assumptions.

A complementary methodological concern is the functional-form sensitivity of parallel trends. \cite{rothsantanna2023parallel} establish that parallel trends is functional-form-invariant if and only if it holds across the entire distribution of untreated potential outcomes, which is a substantially stronger restriction than the unconditional mean version commonly invoked. This insight has direct relevance for absorbing-state outcomes such as project cancellation: where the cumulative cancellation rate mechanically saturates as time progresses, parallel trends in means may fail by construction, motivating the hazard-rate reformulation introduced by \cite{deaner2024causal} and applied in Appendix~\ref{app:deaner_ku} of this thesis.

For settings where the parallel-trends assumption is suspect even under functional-form transformation, \cite{santanna2020doubly} develop doubly-robust estimators that combine outcome modelling with inverse-probability weighting, yielding consistent ATT estimates if either the outcome model or the propensity-score model is correctly specified. The IPWRA estimator applied in Chapter~\ref{ch:results_offtake} for the offtake-effect identification is the doubly-robust class member. \cite{kwonroth2024bayesian} extend the honest-DiD bounds framework to a Bayesian setting, allowing researchers to incorporate prior information about the magnitude or shape of parallel-trends violations rather than relying on data-driven post-period restrictions alone.

The present thesis applies the three primary estimators (TWFE, Sun-Abraham, and BJS-imputation) plus IPWRA for matching settings and assesses convergence as the principal robustness check. When all four estimators agree in sign and approximate magnitude, the substantive conclusion is strengthened.

\subsection{Honest sensitivity analysis}

A separate strand of recent work addresses the limitation that parallel-trends pretest power is low \cite{roth2022pretest}. \cite{rambachan2023more} (RR) propose a framework for \emph{honest} inference: rather than testing parallel trends and conditionally proceeding, the researcher imposes a bound on the magnitude of permissible parallel-trends violations and constructs partial-identification intervals for the average treatment effect.

The RR framework offers two restriction classes. The \emph{relative-magnitudes} restriction (RM) bounds post-treatment trend violations by $M$ times the maximum (or median) observed pre-treatment violation:
\begin{equation}
|\text{bias}^{\text{post}}_t| \leq M \cdot \max_s |\text{bias}^{\text{pre}}_s|.
\label{eq:lit_rr_rm}
\end{equation}
The \emph{smoothness} restriction (SD) bounds the second difference of trends. Both yield identified-set bounds on the ATT.

The breakdown $M^*$ --- the smallest $M$ at which the identified set contains zero --- summarises sensitivity in a single number. Our application of RR (Chapter \ref{ch:results_did}, Section \ref{sec:honest_did}) finds heterogeneous robustness across our four policy comparisons, with substantive interpretive implications.

\subsection{Sensitivity to unobservables}

For cross-sectional causal estimates (as in our offtake-effect analysis, Chapter \ref{ch:results_offtake}), the relevant sensitivity framework is selection-on-unobservables. \cite{oster2019unobservable} extends the \cite{altonji2005selection} approach by jointly bounding bias from two parameters: the coefficient-stability ratio $\delta = \text{cov}(W, \text{unobs}) / \text{cov}(W, \text{obs})$ (how much the unobservables co-vary with treatment relative to observables) and an upper bound $R^2_{\max}$ for the explanatory power of the full set of confounders.

A key diagnostic is the \emph{null-effect} $\delta_{\text{null}}$: the value of $\delta$ at which the treatment effect becomes zero. Conventional thresholds suggest robustness when $|\delta_{\text{null}}| \geq 1$. Our offtake-effect estimate (Chapter \ref{ch:results_offtake}) yields $\delta_{\text{null}} = 20.23$, indicating that unobserved confounders would need to be more than twenty times as strongly correlated with treatment as our extensive observed controls to nullify the effect --- a substantively implausible condition.

\section{Synthesis and positioning}

The present thesis sits at the intersection of three literatures. Methodologically, it inherits modern DiD identification (Sun-Abraham, BJS) and honest sensitivity (Rambachan-Roth, Oster) tools from the causal-inference literature. Theoretically, it extends the real-options literature by linking carrot-policy mechanism design to the Dixit-Pindyck $(V/I, \sigma)$ parameters. Substantively, it provides the first cross-jurisdictional causal evaluation of hydrogen policies, addressing a clear gap in the implementation-gap literature.

The empirical findings (Chapters \ref{ch:results_did}--\ref{ch:results_offtake}) will be shown to converge on the theoretical predictions and to be externally corroborated by the 2024--2026 industry cancellation pattern.
